% Pursuing Stacks (cotes 134-2 to 134-8): archive cote 134-5, site page 107 (PDF page 108), typescript page 294.
% typescript. Draft compared: ps5.tex:148-203.
% AI-transcribed by Claude Opus 5.5 (claude-opus-5-5, adaptive thinking, effort high), Message Batches API, batch ps2 (msgbatch_01NVwr1Yuo8qJf5rykeGBXPK), 2026-11-12; not yet checked by eye.
\scanpage{134-5/107}
\mspage{294}
structure to the corresponding notions of homotopy equivalence between
maps and of homotopisms,  is highly non autodual too.  It breaks down alto-%
gether when $M$ is a ``pointed'' category,  namely \xout{has} contains an object which
is both initial and final \unsure{-}  in this case,  \xout{the} for any homotopy interval
structure on $M$,  any two maps in $M$ are homotopic,  and any map is a homoto-%
pism,  hence any object is contractible !

Our initial motivation,  namely understanding ``modelizers'' for ordi-%
nary homotopy types,  made it very natural to get involved in non autodual
situation,  as the homotopy category (Hot) itself,  and the usual model
categories for it,  displays strongly non autodual features. \hand{(}Thus,  where-%
as the usual test categories all have a final object,  it is easy to see
that a test-category cannot possibly have an initial object.\hand{)} However,
the homotopy and asphericity notions we then came to develop make a sense
and are familiar indeed not only in the ``modelizing story'',  but \xout{practical-}
\xout{ly} in any situation whatever which turns up,  giving rise to anything like
a ``homology'' or a ``homotopy'' theory.  To give one specific example,  star-%
ting with an abelian category,  the corresponding homology theory is con-%
cerned with the category of A-valued \inserted[typed]{\xout{chain}} complexes,  say $K$\hand{$^{\bullet}$(A)}.  \unsure{The} most
basic notions here are the three homotopy notions and two asphericity
notions,  namely:  homotopy equivalence between maps,  homotopisms,  contra\unsure{c}
% NOTE: the end of this line is cut off at the right edge of the scan
tible objects,\inserted{and} quasi-isomorphism\hand{s} (= ``weak equivalence\hand{s}''),  acyclic (=
``aspheric'') objects.  The two first homotopy notions determine each other
mutually in the ususal way,  and define the third,  namely an object is
contractible \xout{iff} \hand{(}or null-homotopic\hand{)} iff the map $0\to X$,  or equivalently $X\to 0$,
is a homotopism.  On the\hand{/}other hand,  if we use the mapping-cylinder constru
% NOTE: the end of this line is cut off at the right edge of the scan
-tion for a map,  the \xout{notion} set of contractible objects determines the set
of homotopisms,  as the maps whose mapping cylinder is contractible.  Like-%
wise,  weak equivalences determine\xout{s} \xout{the notion of} aspheric objects,  and con-%
versely if the mapping cylinder construction is given.  The question that
now comes to mind immediately\inserted{is} \deleted{wather}\inserted{whether} the two sets of notions,  the\inserted{three} pure
homotopy notions (determining each other)\hand{,} and the\inserted{two} ``asphericity notions''
(determining each other too),  mutually determine each other,  as in the
non commutative set-up we have \over{w}{w}orked in so far.  We can also remark that
the functor
\[ H^{0} : K^{\hand{$\bullet$}}(A) \longrightarrow A \]
visibly plays the part of the functor \hand{$\pi$}$_{0}$ in the non commutative set up,
it gives rise moreover to the $H^{i}$ functors (any inte\xout{\unsure{g}}ger $i$) by composing
with the iterated shift functor  (where the shift of $X$ is just the mapping
\margin{\unsure{PJ}}
\margin{107}
